\documentclass{exam}

\newif\ifA
\newif\ifkey
\newif\ifprintselected   % required by exam_config's \ansbox

%%%%%%%%%%%%%%%%%%%%%%%%%%%%%%%%%%%%%%%%%%%%%%%%%%%%%%
%%% SET THESE IF VARIABLES
%%%%%%%%%%%%%%%%%%%%%%%%%%%%%%%%%%%%%%%%%%%%%%%%%%%%%%
% Lab name -- labs are named by the lesson they belong to, never by a
% sequential lab number (Neil, 2026-10-12): this is the L15 Lab.
\def\labLesson{15}

% is this the key or not?
\ifdefined\iskey
	\ifnum\iskey=1
		\keytrue
	\else
		\keyfalse
	\fi
\else
	%% Pick one manually
	\keyfalse % if this is NOT the key
%	\keytrue     % if this IS the key
\fi

% set up the header and footer - change for every term
\ifdefined\thisTermIn
	\def\thisterm{\thisTermIn}
\else
	\def\thisterm{Spring 2027}
\fi
%
\def\thisclass{ECE 448}

%%%%%%%%%%%%%%%%%%%%%%%%%%%%%%%%%%%%%%%%%%%%%%%%%%%%%%
%%% END SET THESE IF VARIABLES
%%%%%%%%%%%%%%%%%%%%%%%%%%%%%%%%%%%%%%%%%%%%%%%%%%%%%%
% packages and shortcuts
\input{myPackages_gr}
\usepackage[hidelinks, linkcolor=red]{hyperref}
\input{myShortcuts}
\setlength{\parindent}{0pt}
\setlength{\parskip}{6pt}

%% exam class customization
\input{exam_config}
% ECE 448 document fonts.
%
% Body text is Barlow (SIL OFL; the faces and the license are in
% latex/fonts/Barlow/). Barlow is not in TeX Live, so it is loaded as an
% OpenType face through fontspec -- which means the document must be built
% with lualatex, not pdflatex. build_practice.sh and build_lab.sh both do.
%
% Compiling with pdflatex anyway is not an error: the guard below falls back
% to the LaTeX default so the source still compiles under any engine and can
% be syntax-checked. The resulting PDF just is not the one we ship.
%
% Paths are relative to latex/, which is where both build scripts run
% (each does cd "$(dirname "$0")"). Loading by path rather than by installed
% family name means no system font install and no fontconfig, so a fresh
% clone builds a byte-comparable PDF anywhere.
%
% MATH IS DELIBERATELY LEFT IN COMPUTER MODERN. Barlow has no math companion,
% and CM math under sans body text is already the house pairing: the decks set
% words in Source Sans Pro and let MathJax render symbols in its CM-derived
% face. Moving math to a sans as well would mean unicode-math plus a sans math
% font (Fira Math being the plausible one) and re-checking every equation in
% the course for symbol coverage. That is a separate decision, not a side
% effect of this one.

\usepackage{iftex}
\ifLuaTeX
	\usepackage{fontspec}
	\defaultfontfeatures{Ligatures=TeX}
	\setmainfont{Barlow}[
		Path           = fonts/Barlow/,
		Extension      = .ttf,
		UprightFont    = *-Regular,
		ItalicFont     = *-Italic,
		BoldFont       = *-Bold,
		BoldItalicFont = *-BoldItalic,
	]
	\setsansfont{Barlow}[
		Path           = fonts/Barlow/,
		Extension      = .ttf,
		UprightFont    = *-Regular,
		ItalicFont     = *-Italic,
		BoldFont       = *-SemiBold,
		BoldItalicFont = *-BoldItalic,
	]
	% microtype earns its place here specifically because of Barlow. Its
	% glyphs are narrower than Computer Modern's, so justified prose sets
	% tighter and TeX has to squeeze interword space to fit; two paragraphs
	% in the L02/L03 sets went overfull on the first Barlow build. Protrusion
	% and font expansion absorb that -- both are lualatex features, so this
	% belongs inside the guard with everything else.
	\usepackage{microtype}

	% Mono stays Latin Modern: Barlow ships no monospaced face, and verbatim
	% and code need one that actually is fixed-pitch.
\fi
   % Barlow body text (needs lualatex)
\printselectedfalse      % full worked solutions in the key, not answers-only
\CorrectChoiceEmphasis{\color{red}}
\SolutionEmphasis{\color{red}}
\renewcommand\questionlabel{\textbf{\thequestion.}}
\renewcommand{\thepartno}{\alph{partno}}
\renewcommand{\partlabel}{(\thepartno)}

% Fillable table cell: blank in the student copy, the expected value in red on
% the KEY. Fixed-width centered columns keep the blank cells writable.
\newcolumntype{K}[1]{>{\centering\arraybackslash}p{#1}}
\newcommand{\kv}[1]{\ifkey{\color{red}#1}\fi\rule[-0.10in]{0pt}{0.34in}}

% Answers are not in this public file. Every solution, \kv value, and \ansbox
% value is \keysol{<id>}, defined in the private ECE448-faculty repo at
% lab-keys/src/ECE448_Lab_L15_MPSK_solutions.tex; build_lab.sh puts that
% folder on TEXINPUTS (KEYDIR) for the KEY. A KEY build without it stops.
\long\def\defkeysol#1#2{\expandafter\long\expandafter\def\csname keysol@#1\endcsname{#2}}
\def\keysol#1{\ifcsname keysol@#1\endcsname\csname keysol@#1\endcsname\fi}
\ifkey
	\IfFileExists{ECE448_Lab_L15_MPSK_solutions.tex}%
		{\input{ECE448_Lab_L15_MPSK_solutions}}%
		{\errmessage{KEY build needs ECE448_Lab_L15_MPSK_solutions.tex from ECE448-faculty/lab-keys/src (set KEYDIR)}}
\fi

% print the answers if this is the key
\ifkey
	\printanswers % if this IS the key
\else
	\noprintanswers  %if this is NOT the key
\fi

% print the header and footer
% need to print key on top if this is the key
\ifkey
	\firstpageheader{}{\color{red}{ KEY}}{}
	\runningheader{\thisclass\hspace{1pt} -- \thisterm }{L\labLesson\ Lab\hspace{1pt}\color{red}{ KEY}}{\thepage\hspace{1pt} of\hspace{1pt} \numpages}
	\firstpagefooter{}{}{}
	\runningfooter{}{}{}
\else
	\firstpageheader{}{}{}
	\runningheader{\thisclass \hspace{1pt} -- \thisterm }{L\labLesson\ Lab\hspace{1pt}}{\thepage\hspace{1pt} of\hspace{1pt} \numpages}
	\firstpagefooter{}{}{}
	\runningfooter{}{}{}
\fi

\runningheadrule


\begin{document}
\pagestyle{headandfoot}   % exam class

%%%%%%%%%%%%%%%%%%%%%%%%%%%%%%%%%%%%%%%%%%%%%%%%%%%%%%
% title block
%%%%%%%%%%%%%%%%%%%%%%%%%%%%%%%%%%%%%%%%%%%%%%%%%%%%%%

\begin{center}
USAF ACADEMY \\
DEPARTMENT OF ELECTRICAL AND COMPUTER ENGINEERING \\
LESSON \labLesson\ LAB - mPSK \\
\textbf{Lab Packet} \\
\thisterm \\[4mm]
\end{center}

\vspace{-4pt}
\textbf{Name:} \rule[-2pt]{2.6in}{0.4pt} \hfill
\textbf{Section:} \rule[-2pt]{0.9in}{0.4pt} \hfill
\textbf{Date:} \rule[-2pt]{1.1in}{0.4pt}

\vspace{4pt}

\fullwidth{\textbf{Learning Objective 6.4: } I can implement an mPSK modulator and
demodulator (BPSK, QPSK, 8PSK) and recover carrier and symbol timing.}

\vspace{2pt}

This packet is the L\labLesson\ Lab, in three parts, and you turn it in once, complete,
before Lesson 16. Do \textbf{Part 1} on your own before the lab period: run the link
flowgraph and work out the settings you will need at the bench. \textbf{Part 2} is
your own measurements, recorded in class: Questions 3 to 6 on your own laptop, and
Question 7 at a PlutoSDR station with your team. Do \textbf{Part 3} on your own after
the lab: explain what you saw and answer the analysis questions.

\textbf{Grading.} You submit the packet once, we grade it against Objective 6.4 on
the course's 1 / 0 scale, and you may revise it under course policy. Paste screenshots
where a question asks for one.

\textbf{Transmitting.} In class we transmit only into a cable, into a shielded
enclosure, or at very low power in an ISM band, and only when the lesson says to.
Questions 1 to 6 run on your laptop and transmit nothing. Question 7 transmits from a
PlutoSDR's TX port, through a cable with an inline 30 dB SMA attenuator, to the same
Pluto's RX port, at 915 MHz. Connect the cable and attenuator before you run the
flowgraph, and never connect an antenna in this lab. Never set TX attenuation below
30 dB or change the slider's range; the RX input is rated to only a few dBm. We have
run the PlutoSDR flowgraph only in simulation, so Question 7 is its first run on
hardware.

The flowgraphs are on the L15 lesson page, \texttt{L15-psk-tx.grc} and
\texttt{L15-psk-link.grc}, and on the L15 Lab page, \texttt{L15-psk-pluto.grc}. Every
number below assumes their settings: 4,000 symbols per second, RRC pulses with
$\alpha = 0.35$ at 8 samples per symbol, Gray-coded points, a Costas Loop bandwidth of
0.02, a default noise amplitude of 0.5, and the default message.

\begin{questions}

\fullwidth{\large\textbf{Part 1: Pre-Lab} \normalsize(individual, before the lab period)}

%%%%%%%%%%%%%%%%%%%%%%%%%%%%%%%%%%%%%%%%%%%%%%%%%%%
\question \textbf{Run the link.} Open \texttt{L15-psk-link.grc} in GRC and run it
with every slider at its default.

\begin{parts}

	\begin{minipage}[t]{0.58\linewidth}
		\part After about 15 s, record the Received label and the Lines label.
		\begin{solution}[0.7in]\keysol{sol1}\end{solution}
	\end{minipage}\hfill%
	\begin{minipage}[t]{0.37\linewidth}
		\raggedleft \vspace{12pt}
		\ansbox[1.4in]{\keysol{ans1}}
	\end{minipage}

	\part Type 50 into the frequency offset box, wait about 5 s, and paste a screenshot of
	the running flowgraph's window.
		\begin{solution}[1.6in]\keysol{sol2}\end{solution}

\end{parts}

%%%%%%%%%%%%%%%%%%%%%%%%%%%%%%%%%%%%%%%%%%%%%%%%%%%
\newpage
\question \textbf{Settings.} You will need each of these at the bench.

\begin{parts}


	\part Question 3 builds an 8PSK flowgraph whose Chunks to Symbols block needs the
	eight points on a circle of radius 1. Using the lesson's 8PSK map, list the angle
	of each point in degrees, by phase index.

	\begin{center}
	\renewcommand{\arraystretch}{1.25}
	\small\setlength{\tabcolsep}{3pt}
	\begin{tabular}{| l | *{8}{K{0.55in} |}}
	\hline
	Index & 0 & 1 & 2 & 3 & 4 & 5 & 6 & 7 \\\hline
	Angle & \kv{\keysol{kv1}} & \kv{\keysol{kv2}} & \kv{\keysol{kv3}} & \kv{\keysol{kv4}} & \kv{\keysol{kv5}} & \kv{\keysol{kv6}} & \kv{\keysol{kv7}} & \kv{\keysol{kv8}} \\\hline
	\end{tabular}
	\end{center}
		\begin{solution}[0.4in]\keysol{sol3}\end{solution}

	\begin{minipage}[t]{0.58\linewidth}
		\part The Costas Loop's frequency output is in radians per symbol. Find the constant
		of the Multiply Const that turns it into hertz at 4,000 symbols per second.
		\begin{solution}[0.9in]\keysol{sol4}\end{solution}
	\end{minipage}\hfill%
	\begin{minipage}[t]{0.37\linewidth}
		\raggedleft \vspace{12pt}
		\ansbox[1.0in]{\keysol{ans2}}
	\end{minipage}

	\begin{minipage}[t]{0.58\linewidth}
		\part Your Question 3 flowgraph adds a 20 Hz frequency offset at 4,000 symbols per
		second. How far does it turn the constellation each symbol, and how many symbols
		does one full turn take?
		\begin{solution}[0.9in]\keysol{sol5}\end{solution}
	\end{minipage}\hfill%
	\begin{minipage}[t]{0.37\linewidth}
		\raggedleft \vspace{12pt}
		\ansbox[1.0in]{\keysol{ans3}}

		\vspace{4pt}
		\ansbox[1.0in]{\keysol{ans4}}
	\end{minipage}

\end{parts}

%%%%%%%%%%%%%%%%%%%%%%%%%%%%%%%%%%%%%%%%%%%%%%%%%%%
\newpage
\fullwidth{\large\textbf{Part 2: Measurements} \normalsize(in class)}

\question \textbf{Build a carrier-recovery loop.} Start GRC with an empty canvas. Do not
open the lesson's flowgraphs: build the 8PSK carrier-recovery flowgraph yourself as the
lab page lists, with a 20 Hz frequency offset.

\begin{parts}

	\part Paste a screenshot of your flowgraph on the GRC canvas. Every connection
	must be black, not red.
		\begin{solution}[1.5in]\keysol{sol6}\end{solution}

	\begin{minipage}[t]{0.58\linewidth}
		\part Paste a screenshot of the running window, and record the loop's frequency
		reading at offsets of 20 Hz and $-60$ Hz.
		\begin{solution}[1.2in]\keysol{sol7}\end{solution}
	\end{minipage}\hfill%
	\begin{minipage}[t]{0.37\linewidth}
		\raggedleft \vspace{12pt}
		20 Hz: \ansbox[1.0in]{\keysol{ans5}}

		\vspace{4pt}
		$-60$ Hz: \ansbox[1.0in]{\keysol{ans6}}
	\end{minipage}

	\part Change the Costas Loop's order from 8 to 2 and run it again. Describe the
	output and the frequency reading, and explain why the loop needs the right order.
		\begin{solution}[1.2in]\keysol{sol8}\end{solution}

\end{parts}

%%%%%%%%%%%%%%%%%%%%%%%%%%%%%%%%%%%%%%%%%%%%%%%%%%%
\newpage
\question \textbf{Noise.} Follow the lab page's steps for \texttt{L15-psk-link.grc}, with
\texttt{diff} = 1, no frequency offset, and the channel phase at 0. For each
constellation, raise the noise amplitude by 0.1 at a time from 0.5, wait about 20 s at
each step, and watch the Costas Loop's output and the Lines label.

\begin{parts}

	\part Record the noise amplitude at which about half the lines received are correct,
	and describe the Costas output there.

	\begin{center}
	\renewcommand{\arraystretch}{1.25}
	\begin{tabular}{| l | K{1.1in} | K{1.1in} | K{1.1in} |}
	\hline
	& BPSK & QPSK & 8PSK \\\hline
	Noise amplitude & \kv{\keysol{kv9}} & \kv{\keysol{kv10}} & \kv{\keysol{kv11}} \\\hline
	\end{tabular}
	\end{center}
		\begin{solution}[1.0in]\keysol{sol9}\end{solution}

	\part At each constellation's limit, describe the Lines label and the Received label:
	how many lines arrive, and what the wrong characters look like. Do the three limits look
	different?
		\begin{solution}[1.2in]\keysol{sol10}\end{solution}

	\part Paste a screenshot of the window for 8PSK at its limit, and mark one wrong
	character in the Received label.
		\begin{solution}[1.4in]\keysol{sol11}\end{solution}

\end{parts}

%%%%%%%%%%%%%%%%%%%%%%%%%%%%%%%%%%%%%%%%%%%%%%%%%%%
\newpage
\question \textbf{Carrier recovery.} Follow the lab page's steps at the default noise
amplitude of 0.5, with differential coding on. Before each offset, start from a loop locked at 0 Hz,
as the lab page describes, and then type the new offset in one jump. Watch the frequency
estimate for up to 20 s.

\begin{parts}

	\part For QPSK, record whether the loop locks, about how long it takes, and the
	frequency it settles at.

	\begin{center}
	\renewcommand{\arraystretch}{1.25}
	\begin{tabular}{| l | K{1.1in} | K{1.1in} | K{1.1in} |}
	\hline
	Offset (Hz) & 250 & 350 & 400 \\\hline
	Locks? & \kv{\keysol{kv12}} & \kv{\keysol{kv13}} & \kv{\keysol{kv14}} \\\hline
	Time & \kv{\keysol{kv15}} & \kv{\keysol{kv16}} & \kv{\keysol{kv17}} \\\hline
	Settles at (Hz) & \kv{\keysol{kv18}} & \kv{\keysol{kv19}} & \kv{\keysol{kv20}} \\\hline
	\end{tabular}
	\end{center}
		\begin{solution}[0.6in]\keysol{sol12}\end{solution}

	\part Set \texttt{M} to 8 and run the flowgraph again. Record the same for 8PSK, and
	find its pull-in range.

	\begin{center}
	\renewcommand{\arraystretch}{1.25}
	\begin{tabular}{| l | K{0.9in} | K{0.9in} | K{0.9in} | K{0.9in} |}
	\hline
	Offset (Hz) & 100 & 150 & 200 & 250 \\\hline
	Locks? & \kv{\keysol{kv21}} & \kv{\keysol{kv22}} & \kv{\keysol{kv23}} & \kv{\keysol{kv24}} \\\hline
	Time & \kv{\keysol{kv25}} & \kv{\keysol{kv26}} & \kv{\keysol{kv27}} & \kv{\keysol{kv28}} \\\hline
	\end{tabular}
	\end{center}
		\begin{solution}[0.6in]\keysol{sol13}\end{solution}

\end{parts}

%%%%%%%%%%%%%%%%%%%%%%%%%%%%%%%%%%%%%%%%%%%%%%%%%%%
\newpage
\question \textbf{Phase ambiguity.} Follow the lab page's steps: set \texttt{diff} to 0 and run
the flowgraph again, so that the loop locks with no frequency offset and the channel phase
at 0, and leave the noise amplitude at 0.5. For each channel phase, start from a loop locked at
0 degrees, as the lab page describes, and then type the phase in one jump. Mark each setting at which
the Lines count keeps rising.

\begin{parts}

	\part Mark the settings that deliver text, and count them.

	\begin{center}
	\renewcommand{\arraystretch}{1.25}
	\footnotesize\setlength{\tabcolsep}{2pt}
	\begin{tabular}{| l | *{18}{K{0.2in} |}}
	\hline
	Phase & 0 & 20 & 40 & 60 & 80 & 100 & 120 & 140 & 160 & 180 & 200 & 220 & 240 & 260 & 280 & 300 & 320 & 340 \\\hline
	QPSK & \kv{\keysol{ph1}} & \kv{\keysol{ph2}} & \kv{\keysol{ph3}} & \kv{\keysol{ph4}} & \kv{\keysol{ph5}} & \kv{\keysol{ph6}} & \kv{\keysol{ph7}} & \kv{\keysol{ph8}} & \kv{\keysol{ph9}} & \kv{\keysol{ph10}} & \kv{\keysol{ph11}} & \kv{\keysol{ph12}} & \kv{\keysol{ph13}} & \kv{\keysol{ph14}} & \kv{\keysol{ph15}} & \kv{\keysol{ph16}} & \kv{\keysol{ph17}} & \kv{\keysol{ph18}} \\\hline
	8PSK & \kv{\keysol{ph19}} & \kv{\keysol{ph20}} & \kv{\keysol{ph21}} & \kv{\keysol{ph22}} & \kv{\keysol{ph23}} & \kv{\keysol{ph24}} & \kv{\keysol{ph25}} & \kv{\keysol{ph26}} & \kv{\keysol{ph27}} & \kv{\keysol{ph28}} & \kv{\keysol{ph29}} & \kv{\keysol{ph30}} & \kv{\keysol{ph31}} & \kv{\keysol{ph32}} & \kv{\keysol{ph33}} & \kv{\keysol{ph34}} & \kv{\keysol{ph35}} & \kv{\keysol{ph36}} \\\hline
	\end{tabular}
	\end{center}
		\begin{solution}[0.5in]\keysol{sol14}\end{solution}

	\part Set \texttt{diff} to 1, run the flowgraph again, and repeat three settings that failed. Record
	what happens. Then, for QPSK with differential coding off, explain what changes if you
	step the phase by 20 degrees at a time without returning to 0.
		\begin{solution}[1.3in]\keysol{sol15}\end{solution}

\end{parts}

%%%%%%%%%%%%%%%%%%%%%%%%%%%%%%%%%%%%%%%%%%%%%%%%%%%
\newpage
\question \textbf{QPSK on a PlutoSDR.} At a PlutoSDR station, with your team, follow the
lab page's steps: connect the Pluto's TX port to its RX port through the 30 dB
attenuator, check the Pluto with \texttt{iio\_info -s}, and run
\texttt{L15-psk-pluto.grc}. When you finish, stop the flowgraph and unplug the Pluto
before you disconnect the cable.

\begin{parts}

	\part Once lines arrive, paste a screenshot that shows both constellations, the
	frequency estimate, and the Lines label.
		\begin{solution}[1.5in]\keysol{sol16}\end{solution}

	\part With \texttt{diff} = 1, stop and start the flowgraph three times. Then set
	\texttt{diff} to 0 and do it eight more times. Record which runs deliver text.

	\begin{center}
	\renewcommand{\arraystretch}{1.25}
	\small\setlength{\tabcolsep}{3pt}
	\begin{tabular}{| l | *{3}{K{0.27in} |} *{8}{K{0.27in} |}}
	\hline
	Run & 1 & 2 & 3 & 1 & 2 & 3 & 4 & 5 & 6 & 7 & 8 \\\hline
	Differential & on & on & on & off & off & off & off & off & off & off & off \\\hline
	Text? & & & & & & & & & & & \\\hline
	\end{tabular}
	\end{center}
		\begin{solution}[0.8in]\keysol{sol17}\end{solution}

	\part With \texttt{diff} = 1, wait 30 s at each TX attenuation and record the Lines label.
	Explain what changes as the attenuation rises.

	\begin{center}
	\renewcommand{\arraystretch}{1.25}
	\begin{tabular}{| l | K{1.1in} | K{1.1in} | K{1.1in} |}
	\hline
	TX attenuation (dB) & 70 & 80 & 89 \\\hline
	Lines correct of received & & & \\\hline
	\end{tabular}
	\end{center}
		\begin{solution}[1.2in]\keysol{sol18}\end{solution}

\end{parts}

%%%%%%%%%%%%%%%%%%%%%%%%%%%%%%%%%%%%%%%%%%%%%%%%%%%
\newpage
\fullwidth{\large\textbf{Part 3: Writeup} \normalsize(individual, after the lab)}

\question \textbf{Analysis.}

\begin{parts}

	\part Compare the noise limits in Question 4 for the three constellations. Explain
	in words why they come in that order, and what that means for a radio that must carry
	more bits per symbol.
		\begin{solution}[1.5in]\keysol{sol19}\end{solution}

	\part Compare the 8PSK pull-in range with QPSK's, explain the difference, and say what a link between two different Plutos,
	each accurate to 25 ppm at 915 MHz, would need before its Costas Loop.
		\begin{solution}[1.5in]\keysol{sol20}\end{solution}

	\part Compare Question 6 with Question 7(b). Explain why the
	working settings were where they were, why the Pluto runs behaved as they did, and
	why a real mPSK link uses differential coding even though one wrong symbol then
	corrupts two.
		\begin{solution}[1.5in]\keysol{sol21}\end{solution}

\end{parts}

\end{questions}

\vspace{1.0em}
\noindent\textbf{Documentation:} \rule[-2pt]{0.6\linewidth}{0.4pt}

\end{document}
